\scanpage{108}
\uline{Théorème} $G^{\unsure{0|\circ}}$ est transitif sur l'ens des arcs de $(S,A)$ (\unsure{lire|i.e.} des couples $(s,t)$ de deux pts liés) et a fortiori sur l'ens des sommets, \deleted{et} \inserted{l'ens des arêtes et} sur l'ens des triangles (\inserted{ensemble} qui sont isom. à des ens. quotients de l'ens précédent.)

% slanted
\margin{Sauf exception $c\,(=\operatorname{Card} E_a)=1$ i.e. $(S,A)$ est le bi-triangle. Alors $G$ a les transitivités \unsure{énoncées} pour $G^0\inserted{\unsure{\simeq}\mathfrak S_3\times\mathfrak S_3}$, qui est d'indice $2$ dans $G$ (\deleted{\unsure{Aut}} $\mathfrak S_2.(\mathfrak S_3\times\mathfrak S_3)$). $G^0$ est encore transitif sur $S$, mais pas sur $\deleted{\unsure{A}}T$ (ni à fortiori sur $A$) -- il est transitif sur chacune des deux familles de triangles parallèles\ldots}

\uline{Cor} Les $\sigma_{\underline t}$ $(\underline t\in T)$ sont conjugués \inserted{dans $G^0$ (a fortiori en $G$)}, et forment exactement une classe de conjugaison d'éléments de $G^0$ (ou de $G$)

Résulte des deux lemmes suivants

\uline{Lemme 1} Soient $s,t\in S$, alors

a) Si $s,t$ sont liés, $\exists\ \underline t\in T$ tel que $t=\sigma_{\underline t}s$

b) \deleted{Si $s,t$ sont non liés} \inserted{Dans tous les cas} $\exists\ u\in S$ tel que $s$ lié à $u$ et $u$ lié à $t$, donc (vu a) des $\underline t,\underline t'\in T$ tels que $u=\sigma_{\underline t}s$, $t=\sigma_{\underline t'}u$, donc $t=\sigma_{\underline t'}\sigma_{\underline t}s$.

Le lemme a) est clair (\deleted{prendre pour $\underline t$} \inserted{soit $a$ le troisième sommet du triangle défini par $\{s,t\}$}, \deleted{le triangle défini \unsure{ces trois|un triangle}} et soit $\underline t$ \unsure{un} triangle de sommet $a$ distinct du précédent) ; d'où immédiatement b.

\uline{Lemme 2} \inserted{(\unsure{Supposons|Supposant} $c\neq1$, i.e. il y a au moins $3$ triangles incidents en chaque sommet\ldots} Soient $a\in S$ et $s,t$ deux él. de $S$ liés à $a$, \inserted{$\unsure{s\neq t}$ et $s$ non lié à $t$.} Alors $\exists\,g\in G^0$ (produit \deleted{de deux} $\sigma_{\unsure{\underline\tau|\underline t}} \sigma_{\underline t}$, où $\underline\tau,\underline t$ sont $2$ triangles distincts convenables) tels que $u(a)=a$, $u(s)=t$, $u(t)=s$

% archive sheet number: 107
